% Part: lambda-calculus
% Chapter: introduction
% Section: fixed-point-combinator

\documentclass[../../../include/open-logic-section]{subfiles}

\begin{document}

\olfileid{lam}{int}{fix}
\olsection{Kombinator Titik Tetep}

Upamana ana suku lambda $g$, lan kita
ngarepake suku liyane $k$ kang nduweni
sipat yen $k$ ekuivalen-$\beta$ karo
$gk$. Tetepna suku
\[
\fn{diag}(x) = xx
\]
lan
\[
l(x) = g(\fn{diag}(x))
\]
miturut konvensi notasi kita; tegese,
$l$ iku suku $\lambd[x][g(xx)]$.
Upamana $k$ iku suku $ll$. Banjur
kita duwe
\begin{align*}
k & = (\lambd[x][g(xx)])(\lambd[x][g(xx)]) \\
& \red  g((\lambd[x][g(xx)])(\lambd[x][g(xx)])) \\
& = gk.
\end{align*}
Yen kita njupuk
\[
Y = \lambd[g][((\lambd[x][g(xx)])(\lambd[x][g(xx)]))]
\]
mula $Yg$ lan $g(Yg)$ loro-lorone bisa
direduksi dadi suku kang padha; dadi
$Yg \equiv_\beta
g(Yg)$. Iki diarani ``kombinator Curry.''
Yen kita njupuk
\[
Y = (\lambd[xg][g(xxg)])(\lambd[xg][g(xxg)])
\]
mula $Yg$ malah direduksi dadi
$g(Yg)$, pratelan kang luwih kuwat.
Wujud $Y$ kang kapindho iki diarani
``kombinator Turing.''

\end{document}
